\section{Activation Capping}\label{sec:activation_capping}

We use an approach to activation capping motivated by the approach in \citet{lu2026assistant}. For each OCEAN LoRA, we precompute a per-layer activation \emph{axis}
that captures the direction in the residual stream along which the LoRA
shifts the model's hidden states. Activation capping at inference time
clamps the projection of the live residual onto this axis to a fixed
fraction of the LoRA-induced range

We use 240 open-ended prompts from (generated by Claude Opus 4.7). The prompts are intentionally not specific to any OCEAN trait
so that the captured axis reflects the generic shift induced by the
persona LoRA rather than topic-driven differences.

For each prompt, we generate $N=5$ rollouts at
temperature $1.0$, max $256$ new tokens, with two model configurations:
the unmodified base (Llama-3.1-8B-Instruct) and the same model with the
persona LoRA applied. Both sets of rollouts use the same prompts and seed
so the only difference is the applied LoRA.

We then re-run each generated
(prompt, response) pair through both models and extract the per-token
residual-stream activations for the \emph{response} tokens at every
layer. For each layer we take the mean activation vector across all tokens and responses for each group, then the activation direction is taken as the vector pointing from the centroid from the responses from the model without the LoRA, to the centroid from the responses from the model with the LoRA. This defines our axis that we will use for capping.

We define the minimum/maximum value (for a given layer) as the minimum/maximum projection onto this vector for any given response (averaged over all tokens) from either the model with or without the LoRA.

We apply vector capping to the final 40\% of layers. At inference time, for each token, we compute the current projection of the residual stream onto the capping axis and compare it to a target value $proj_{capped}$. We then add a vector along the capping axis (parallel or anti-parallel as needed) whose magnitude exactly closes the gap, leaving the residual unchanged in all directions orthogonal to the axis. The target projection is defined as:
\[
  proj_{capped} \;=\;
  \begin{cases}
    \max(proj, vec\_lim), & vec\_lim \geq 0\\
    \min(proj, 1 - vec\_lim), & vec\_lim < 0
  \end{cases}
\]

So when $vec\_lim$ is positive it creates a floor, and when it is negative it creates a ceiling. The floor can go above 1.0 and the ceiling can go below 0.0, essentially becoming steering vectors. This dynamic approach was chosen to allow us to do sweeps that could be compared to our LoRA scale sweeps.

See sweeps for our OCEAN LoRAs \Cref{fig:actcap_amplifiers,fig:actcap_suppressors}. Whilst the \texttt{TRAIT} scores are relatively flat, suggesting that this intervention doesn't make the model self identify differently, the LLM judge scores do clearly show the expected directional change in each OCEAN trait.



\newcommand{\actcapfigdir}{figures/appendix/activation_capping_mcq_llm_judge_evals}

\newcommand{\actrow}[2]{%
    \includegraphics[width=0.27\linewidth]{\actcapfigdir/trait_sweep_#1_#2_actcap.pdf}\hfill
    \includegraphics[width=0.32\linewidth]{\actcapfigdir/mmlu_breakdown_#1_#2_actcap.pdf}\hfill
    \includegraphics[width=0.32\linewidth]{\actcapfigdir/judge_sweep_#1_#2_actcap.pdf}\\[4pt]
}

% Generated by: scripts/visualisations/appendix_actcap_{trait,mmlu,judge}.py
% Data: persona-cartography/monorepo fine_tuning/llama-3.1-8b-it/ocean/*/ocean_const_paired_dpo/evals/{mcq/activation_capping/, llm_judge_activation_capping_sweep/}
\begin{figure}[htbp]
    \centering
    \actrow{openness}{plus}
    \actrow{conscientiousness}{plus}
    \actrow{extraversion}{plus}
    \actrow{agreeableness}{plus}
    \actrow{neuroticism}{plus}
    \caption{Activation capping sweeps for the OCEAN amplifiers. Rows are O$\uparrow$, C$\uparrow$, E$\uparrow$, A$\uparrow$, N$\uparrow$ (top to bottom). Columns are TRAIT logprob sweep (left), MMLU breakdown (middle), Qwen3-235B LLM-judge sweep with coherence on the secondary axis (right). A scale of $0$ corresponds to no activation capping (the unmodified base+adapter model). All error bars are 95\% confidence intervals: BCa bootstrap (1000 resamples) for the trait logprob and judge/coherence scores, Wilson score interval for the MMLU per-category fractions.}
    \label{fig:actcap_amplifiers}
\end{figure}

\begin{figure}[htbp]
    \centering
    \actrow{openness}{minus}
    \actrow{conscientiousness}{minus}
    \actrow{extraversion}{minus}
    \actrow{agreeableness}{minus}
    \actrow{neuroticism}{minus}
    \caption{Activation capping sweeps for the OCEAN suppressors. Rows are O$\downarrow$, C$\downarrow$, E$\downarrow$, A$\downarrow$, N$\downarrow$ (top to bottom). Columns are TRAIT logprob sweep (left), MMLU breakdown (middle), Qwen3-235B LLM-judge sweep with coherence on the secondary axis (right). A scale of $0$ corresponds to no activation capping (the unmodified base+adapter model). All error bars are 95\% confidence intervals: BCa bootstrap (1000 resamples) for the trait logprob and judge/coherence scores, Wilson score interval for the MMLU per-category fractions.}
    \label{fig:actcap_suppressors}
\end{figure}
\FloatBarrier